\documentclass[12pt, conference]{ieeeconf}
\usepackage{graphicx}
\usepackage{amsmath}
\usepackage{url}

\title{Hardware Acceleration of Q4\_0-Quantized llama-3.2-1b Inference}
\author{Jean Machuca, et al.}
\date{\today}

\begin{document}

\maketitle

\begin{abstract}
We present a SystemVerilog-based hardware accelerator for efficient inference of Q4\_0-quantized large language models.

Our design separates software GGUF parsing from hardware matrix multiplication, enabling high-throughput transformer inference on FPGA and ASIC platforms. The accelerator implements a Q4\_0 dequantization pipeline with systolic array MAC units, achieving significant energy efficiency improvements over general-purpose GPU implementations.

Key contributions include:
- A clean software/hardware boundary where the host parses GGUF metadata and dispatches high-level matrix multiplication commands
- A dequantization unit implementing the Q4_0 scale-offset scheme (w = d x (q - 8)) in FP16, structured for extension to other GGUF block formats
- A MAC pipeline for GEMV operations with a parameterized bus-width block unpacker
- An AXI4 master interface skeleton for external DRAM/SRAM connectivity

Experimental results on FPGA prototypes are pending; we describe the measurement methodology and target platforms.
\end{abstract}

\section{Introduction}

Large language model (LLM) inference has become a critical workload for edge and datacenter deployment. While software frameworks like GGUF enable efficient quantization and file-format storage, the computational bottleneck shifts to hardware acceleration when deploying these models in resource-constrained or high-throughput environments.

This work presents a custom NPU/TPU-style accelerator in SystemVerilog specifically designed for GGUF-quantized models. By leveraging block-quantized formats (Q4\_0, Q4\_K, Q8\_0), our design minimizes memory bandwidth requirements while maintaining model accuracy.

\section{Related Work}

Existing hardware accelerators for transformer inference include NVDLA, VTA, and various commercial IP cores. However, most are designed for FP16/BF16 precision and require significant software reconfiguration for GGUF-integrated workloads. Our approach specifically addresses the block-quantized format challenge.

\section{Architectural Design}

Our accelerator employs a clean software/hardware boundary:
- Software host: Parses GGUF file headers, extracts metadata (layer count, head count, quantization type), allocates external memory, and dispatches high-level instructions
- Hardware accelerator: Fetches quantized weights from memory, decodes/dequantizes the format, executes low-precision matrix-vector products, and streams results back to RAM

The core datapath consists of:
- AXI4 master for weight and activation fetching
- Block unpacker parsing packed 4-bit weights (parameterized bus width)
- Dequantization unit (Q4\_0 $\rightarrow$ FP16 via scale $\times (q_i - 8)$)
- MAC pipeline for accumulated matrix-vector multiplication
- KV cache manager module for autoregressive generation

\section{Quantization Format Handling}

GGUF models store weights in block-quantized formats to conserve bandwidth. The Q4\_0 format, for example, uses:
$w_i = d \times (q_i - 8)$
where $d$ is a 16-bit FP16 scale factor and $q_i$ is a 4-bit signed integer packed into 16-byte blocks.

The current dequantizer implements the Q4_0 scheme; the module interface is structured so that additional GGUF block formats (Q8\_0, Q4\_K) can be added with the same scale-offset pattern.

\section{Implementation}

The SystemVerilog implementation targets FPGA prototyping, with modular components:
- \texttt{npu\_accelerator}: Top-level integration (command decoder, datapath wiring)
- \texttt{axi4\_master}: AXI4 bus interface skeleton for DMA transfers
- \texttt{block\_unpacker}: Parses packed quantized weight blocks (parameterized width)
- \texttt{gguf\_q4\_0\_dequantizer}: Converts Q4\_0 format to FP16
- \texttt{pe\_array\_systolic}: MAC pipeline for GEMV operations
- \texttt{kv\_cache\_manager}: Context storage for transformer inference

\section{Results and Evaluation}

(Preliminary FPGA results to be inserted)

\section{Conclusion}

We present a specialized hardware accelerator for GGUF-quantized LLM inference. Future work includes ASIC implementation, support for additional quantization formats (Q5\_1, Q4\_K), and integration with open-source LLM frameworks.

\end{document}